% !TEX program = pdflatex
% Praca dyplomowa -- układ zgodny z szablonem (strona tytułowa i strona AI z szablonu Overleaf
% „Praca dyplomowa AGH”, reszta wg wzoru pracy z Katedry Telekomunikacji WIET).
% Do wgrania do projektu (z szablonu Overleaf): logo_AGH.pdf, AI_poziom_0.jpg ... AI_poziom_3.jpg
\documentclass[12pt,a4paper]{report}

\usepackage{polski}
\usepackage[utf8]{inputenc}
\usepackage[T1]{fontenc}
\usepackage{mathptmx}                 % Times, jak we wzorze pracy
\usepackage[a4paper,left=3.5cm,right=2.5cm,top=2.5cm,bottom=2.5cm]{geometry}
\usepackage{setspace}
\onehalfspacing
\usepackage{indentfirst}
\setlength{\parindent}{1.25cm}
\usepackage{xcolor}
\usepackage{graphicx}
\usepackage{booktabs}
\usepackage{amsmath,amssymb}
\usepackage{float}
\usepackage{enumitem}
\setlist{leftmargin=2cm,topsep=2pt,itemsep=0pt}
\usepackage{chngcntr}
\counterwithout{equation}{chapter}    % numeracja wzorów ciągła: (1), (2), ...
\usepackage{csquotes}
\usepackage{listings}
\lstset{basicstyle=\ttfamily\small,breaklines=true,frame=single,extendedchars=true,literate={ś}{{\'s}}1 {ó}{{\'o}}1 {ł}{{\l}}1 {ą}{{\k a}}1 {ę}{{\k e}}1 {ż}{{\.z}}1 {ź}{{\'z}}1 {ć}{{\'c}}1 {ń}{{\'n}}1}
\usepackage{titlesec}
\usepackage{tocloft}
\usepackage{fancyhdr}
\usepackage[backend=biber,style=numeric,sorting=none]{biblatex}
\addbibresource{bib/literatura.bib}
\usepackage[hidelinks]{hyperref}

% --- nagłówki jak we wzorze: „1. Tytuł” i „1.1 Tytuł” (pogrubiona kursywa)
\titleformat{\chapter}{\normalfont\Large\bfseries}{\thechapter.}{0.5em}{}
\titlespacing*{\chapter}{0pt}{30pt}{28pt}
\titleformat{\section}{\normalfont\large\bfseries\itshape}{\thesection}{0.6em}{}
\titlespacing*{\section}{1.25cm}{14pt}{6pt}
\titleformat{\subsection}{\normalfont\normalsize\bfseries\itshape}{\thesubsection}{0.6em}{}
\titlespacing*{\subsection}{1.25cm}{10pt}{4pt}

% --- spis treści z kropkami
\renewcommand{\cftchapfont}{\normalfont\bfseries}
\renewcommand{\cftchapleader}{\cftdotfill{\cftdotsep}}
\renewcommand{\cftchappagefont}{\normalfont}

% --- numer strony w stopce: „- 13 -”
\pagestyle{fancy}\fancyhf{}\fancyfoot[C]{- \thepage{} -}\renewcommand{\headrulewidth}{0pt}
\fancypagestyle{plain}{\fancyhf{}\fancyfoot[C]{- \thepage{} -}\renewcommand{\headrulewidth}{0pt}}

% --- podpisy: „Rys. 1.1 Opis” + linia „Źródło: [n]”
\usepackage{caption}
\AtBeginDocument{\renewcommand{\figurename}{Rys.}\renewcommand{\tablename}{Tab.}}
\captionsetup{labelfont=bf,labelsep=space,justification=centering}
\newcommand{\zrodlo}[1]{{\centering Źródło: #1\par}}

\newcommand{\todo}[1]{\textcolor{red}{[TODO: #1]}}

% --- dane do strony tytułowej (uzupełnij)
\newcommand{\katedra}{KATEDRA TELEKOMUNIKACJI}      % \todo{katedra opiekuna pracy}
\newcommand{\autor}{\todo{Imię i Nazwisko}}
\newcommand{\opiekun}{\todo{stopień/tytuł Imię i Nazwisko}}
\newcommand{\rokpracy}{2026}
% Poziom użycia AI (0--3), wybierz zgodnie z faktycznym użyciem i wytycznymi wydziału
\newcommand{\poziomAI}{2}

\begin{document}
\newgeometry{margin=3.5cm}
\thispagestyle{empty}
\begin{center}
\IfFileExists{logo_AGH.pdf}{\includegraphics[width=\textwidth]{logo_AGH.pdf}}{\fbox{\parbox{0.8\textwidth}{\centering\vspace{1.5cm}wgraj logo\_AGH.pdf z szablonu Overleaf\vspace{1.5cm}}}}\\[5mm]
{\bf\sf WYDZIAŁ INFORMATYKI, ELEKTRONIKI I TELEKOMUNIKACJI}\\[5mm]
{\bf\sf\katedra}\\[14mm]
{\sf\huge Projekt dyplomowy}\\[12mm]
{\sf\Large Projekt i implementacja webowego systemu detekcji anomalii w transakcjach płatniczych z wykorzystaniem porównawczej analizy algorytmów uczenia maszynowego\\[2mm]
Design and implementation of a web-based anomaly detection system for payment transactions using a comparative analysis of machine learning algorithms}\\[30mm]
\end{center}
{\sf
\begin{tabular}{ll}
Autor: & \autor\\
Kierunek studiów: & Teleinformatyka\\
Opiekun pracy: & \opiekun\\
\end{tabular}
}\\[10mm]
\begin{center}
{\sf Kraków, \rokpracy}
\end{center}
\restoregeometry
\newpage

% Strona z poziomem użycia AI -- zostaje TYLKO JEDNA ikona (ustawiana w main.tex: \poziomAI).
\thispagestyle{empty}
\vspace*{3cm}
{\sf
\begin{tabular}{lp{10cm}}
\IfFileExists{AI_poziom_\poziomAI.jpg}{\includegraphics{AI_poziom_\poziomAI.jpg}}{\fbox{AI\_poziom\_\poziomAI.jpg}} &
\ifcase\poziomAI Brak AI\or Korzystanie tylko do celów redakcyjnych\or Korzystanie tylko po odpowiednim udokumentowaniu\or Korzystanie jest dozwolone bez żadnych ograniczeń\fi\\
\end{tabular}
}
\newpage

% Treść ze wzoru pracy (2018). Powołane akty prawne mogą być nieaktualne -- sprawdź aktualny wzór oświadczenia w dziekanacie/USOS.
\thispagestyle{empty}
\vspace*{2cm}
\noindent Uprzedzony o odpowiedzialności karnej na podstawie art. 115 ust. 1 i 2 ustawy z dnia 4 lutego 1994 r. o prawie autorskim i prawach pokrewnych (t.j. Dz.U. z 2006 r. Nr 90, poz. 631 z późn. zm.): \enquote{Kto przywłaszcza sobie autorstwo albo wprowadza w błąd co do autorstwa całości lub części cudzego utworu albo artystycznego wykonania, podlega grzywnie, karze ograniczenia wolności albo pozbawienia wolności do lat 3. Tej samej karze podlega, kto rozpowszechnia bez podania nazwiska lub pseudonimu twórcy cudzy utwór w wersji oryginalnej albo w postaci opracowania, artystyczne wykonanie albo publicznie zniekształca taki utwór, artystyczne wykonanie, fonogram, wideogram lub nadanie.}, a także uprzedzony o odpowiedzialności dyscyplinarnej na podstawie art. 211 ust. 1 ustawy z dnia 27 lipca 2005 r. Prawo o szkolnictwie wyższym (t.j. Dz. U. z 2012 r. poz. 572, z późn. zm.) \enquote{Za naruszenie przepisów obowiązujących w uczelni oraz za czyny uchybiające godności studenta student ponosi odpowiedzialność dyscyplinarną przed komisją dyscyplinarną albo przed sądem koleżeńskim samorządu studenckiego, zwanym dalej \enquote{sądem koleżeńskim}}, oświadczam, że niniejszą pracę dyplomową wykonałem(-am) osobiście i samodzielnie i że nie korzystałem(-am) ze źródeł innych niż wymienione w pracy.

\vspace{2cm}
\begin{flushright}
………………………
\end{flushright}
\newpage

% \chapter*{Streszczenie}
\addcontentsline{toc}{chapter}{Streszczenie}
\todo{ok. 150--250 słów: problem, metoda, wyniki, wnioski (napisz na końcu własnymi słowami).}

\textbf{Słowa kluczowe:} detekcja anomalii, wykrywanie oszustw, uczenie maszynowe, aplikacja webowa

\chapter*{Abstract}
\addcontentsline{toc}{chapter}{Abstract}
\todo{English version.}

\textbf{Keywords:} anomaly detection, fraud detection, machine learning, web application
   % opcjonalnie, jeśli wydział wymaga streszczenia w pracy

\setcounter{tocdepth}{1}
\tableofcontents

% SZKIC: tekst przygotowany z pomocą AI. Przepisz go własnymi słowami, sprawdź źródła i uwzględnij to w deklaracji poziomu AI.
\chapter*{Wstęp}
\addcontentsline{toc}{chapter}{Wstęp}

Płatności elektroniczne należą do podstawowych usług sektora finansowego, a wraz z ich upowszechnieniem rośnie znaczenie wykrywania nadużyć. Każda transakcja musi być oceniona w krótkim czasie, a oszustwa stanowią znikomą część wszystkich operacji. Wykrywanie oszustw płatniczych jest więc jednocześnie problemem klasyfikacji silnie niezbalansowanej, problemem detekcji anomalii oraz problemem inżynierskim: system oceniający transakcje musi działać szybko, być obserwowalny i pozwalać na zarządzanie modelami \cite{dalpozzolo2018}. \todo{1--2 zdania z liczbami o skali oszustw płatniczych, ze źródłem (np. raport EBA/ECB).}

W literaturze porównuje się zwykle metody uczenia nadzorowanego, które korzystają z etykiet oszustw, oraz nienadzorowanego, które szukają transakcji odbiegających od normy bez etykiet \cite{niu2019,carcillo2019}. Wyniki takich porównań zależą od sposobu podziału danych i doboru miar jakości, a sama ocena modeli rzadko jest połączona z narzędziem, w którym użytkownik mógłby zarządzać modelami i obserwować ich działanie.

Celem pracy jest zaprojektowanie i implementacja webowego systemu do wykrywania anomalii w transakcjach płatniczych oraz porównanie skuteczności algorytmów uczenia nadzorowanego i nienadzorowanego na silnie niezbalansowanych danych. Metodykę porównania oparto na pracy \cite{niu2019}, która jest najbliższa tematycznie, i rozszerzono ją o podział czasowy danych, miary właściwe dla klas niezbalansowanych oraz dobór progu na zbiorze walidacyjnym. Wkład własny stanowią: system webowy z modułowym interfejsem do źródeł transakcji, zarządzaniem modelami i monitorowaniem metryk, odtworzenie i rozszerzenie eksperymentów z literatury w ujednoliconej metodyce oraz analiza wpływu metryki, progu i sposobu radzenia sobie z niezbalansowaniem na wnioski z porównania. \todo{doprecyzować po rozmowie z promotorem}

W pracy postawiono następujące pytania badawcze:
\begin{itemize}
  \item na ile modele nadzorowane przewyższają nienadzorowane (PR-AUC, F1) przy podziale czasowym danych,
  \item jak wybór metryki (ROC AUC a PR-AUC) i progu decyzyjnego wpływa na wnioski z porównania,
  \item jak sposób radzenia sobie z niezbalansowaniem klas (wagi klas, SMOTE) wpływa na modele nadzorowane,
  \item czy system działa w czasie zbliżonym do rzeczywistego (opóźnienie oceny pojedynczej transakcji).
\end{itemize}
Hipoteza robocza zakłada, że modele nadzorowane osiągną wyższy PR-AUC, ale nienadzorowane pozostaną użyteczne, zwłaszcza gdy etykiet brakuje lub są opóźnione.

Eksperymenty przeprowadzono na publicznym zbiorze \emph{Credit Card Fraud Detection} \cite{kaggle_ulb}. Praca ma charakter akademicki: \enquote{czas rzeczywisty} jest symulowany przez odtwarzanie zapisanych transakcji, a wdrożenie produkcyjne i integracja z rzeczywistym systemem płatniczym pozostają poza zakresem. Architekturę systemu zaprojektowano modułowo, tak aby możliwe było dołączenie kolejnych źródeł transakcji (własne pliki, interfejsy bankowe), co opisano jako kierunek dalszego rozwoju. \todo{uaktualnić, jeśli PoC z API bankowym wejdzie do zakresu}

W pierwszym rozdziale opisano domenę, czyli operacje płatnicze i problem oszustw. Rozdział drugi przedstawia metody uczenia maszynowego, miary jakości i przegląd literatury. W rozdziałach trzecim i czwartym zaprezentowano projekt oraz implementację systemu, a w piątym badania własne i analizę wyników. Ostatnią częścią pracy jest podsumowanie, w którym zebrano najważniejsze wnioski.

% SZKIC: tekst przygotowany z pomocą AI. Przepisz go własnymi słowami, sprawdź źródła i uwzględnij to w deklaracji poziomu AI.
\chapter{Operacje płatnicze i problem wykrywania oszustw}
\label{ch:platnosci}

\section{Transakcje płatnicze i systemy płatności}
Transakcja płatnicza jest zleceniem przekazania środków od płatnika do odbiorcy, realizowanym za pośrednictwem instrumentu płatniczego (karta, przelew, portfel mobilny). Z punktu widzenia systemu antyfraudowego istotne są dane opisujące zdarzenie: kwota, czas, odbiorca i jego kategoria, kanał (płatność w sklepie lub online), urządzenie oraz historia klienta. \todo{krótki opis uczestników: płatnik, akceptant, agent rozliczeniowy, wydawca karty; źródło}

\section{Rodzaje oszustw płatniczych i skala problemu}
Najczęściej wyróżnia się nadużycia z użyciem karty bez jej fizycznej obecności (płatności internetowe), wykorzystanie skradzionej lub sklonowanej karty oraz przejęcie konta klienta. \todo{definicje i dane liczbowe z raportu EBA/ECB o oszustwach płatniczych, z cytowaniem} Regulacje europejskie (PSD2) wprowadziły m.in. silne uwierzytelnianie klienta, co wpływa na strukturę oszustw i na dane dostępne dla systemów detekcji. \todo{źródło dla PSD2/SCA}

\section{Systemy antyfraudowe: reguły biznesowe a uczenie maszynowe}
Klasyczne systemy opierają się na regułach definiowanych przez ekspertów (limity kwot, kraje wysokiego ryzyka, częstość transakcji). Są one czytelne, lecz wymagają ręcznej aktualizacji i trudno wykrywają nowe schematy. Uczenie maszynowe pozwala budować modele na historycznych danych i oceniać ryzyko pojedynczej transakcji \cite{dalpozzolo2018}, zazwyczaj uzupełniając reguły, a nie je zastępując.

\section{Specyfika danych transakcyjnych}
Dane transakcyjne mają cechy, które determinują dobór metod i sposób oceny:
\begin{itemize}
  \item \textbf{Silne niezbalansowanie klas.} Oszustwa stanowią ułamek procenta transakcji; w zbiorze użytym w pracy jest to 0{,}17\% (rozdz.~\ref{ch:projekt}).
  \item \textbf{Dryf koncepcji.} Zachowania klientów i sposoby oszustw zmieniają się w czasie, więc model trenowany na przeszłości może tracić skuteczność \cite{dalpozzolo2018}.
  \item \textbf{Opóźnienie weryfikacji.} Etykieta „oszustwo” jest znana dopiero po zgłoszeniu lub dochodzeniu, więc najnowsze dane są nieoznaczone lub błędnie oznaczone \cite{dalpozzolo2018}.
  \item \textbf{Dane wrażliwe.} Publiczne zbiory są anonimizowane, co ogranicza interpretowalność cech.
\end{itemize}
Z dwóch pierwszych cech wynika uzasadnienie podziału czasowego danych i stosowania miar właściwych dla klas niezbalansowanych (rozdz.~\ref{ch:ml}).

\section{Wymagania praktyczne}
System detekcji musi ocenić transakcję w czasie, który nie opóźnia płatności, a decyzja o progu alertu zależy od kosztów błędów: pominięte oszustwo (fałszywie ujemne) oznacza stratę, a fałszywy alarm (fałszywie dodatni) koszt obsługi i pogorszenie doświadczenia klienta. Konieczna jest także możliwość wyjaśnienia decyzji oraz monitorowania jakości modeli w czasie. Wymagania te przekładają się na funkcje opisanego w pracy systemu (rozdz.~\ref{ch:projekt}).

% SZKIC: tekst przygotowany z pomocą AI. Przepisz go własnymi słowami, sprawdź źródła i uwzględnij to w deklaracji poziomu AI.
\chapter{Metody uczenia maszynowego w detekcji anomalii}
\label{ch:ml}

\section{Detekcja anomalii: ujęcie ogólne}
Anomalia to obserwacja istotnie odbiegająca od zachowania uznawanego za normalne. Wyróżnia się anomalie punktowe, kontekstowe (nietypowe w danym kontekście, np. o danej porze) i zbiorowe \cite{chandola2009}. Wykrywanie oszustw płatniczych można ujmować na dwa sposoby: jako klasyfikację binarną na danych z etykietami (uczenie nadzorowane) albo jako szukanie obserwacji odstających bez użycia etykiet (uczenie nienadzorowane). Oba podejścia porównuje się w literaturze \cite{niu2019} i łączy w rozwiązaniach hybrydowych \cite{carcillo2019}.

\section{Algorytmy uczenia nadzorowanego}
\subsection{Regresja logistyczna}
Model liniowy zwracający prawdopodobieństwo klasy dodatniej $p(y{=}1\mid\mathbf{x})=\sigma(\mathbf{w}^\top\mathbf{x}+b)$, gdzie $\sigma(z)=1/(1+e^{-z})$. Parametry wyznacza się minimalizując entropię krzyżową z regularyzacją. Służy jako punkt odniesienia (ang. \emph{baseline}).

\subsection{Drzewa decyzyjne i Random Forest}
Drzewo decyzyjne dzieli przestrzeń cech sekwencją podziałów wybieranych tak, aby zmniejszać niejednorodność węzłów, np. nieczystość Giniego $G=1-\sum_k p_k^2$ \cite{breiman1984}. Pojedyncze drzewa łatwo się przeuczają. Random Forest buduje wiele drzew na próbach bootstrapowych, losując przy każdym podziale podzbiór cech, a decyzję podejmuje przez uśrednienie lub głosowanie \cite{breiman2001}, co zmniejsza wariancję modelu.

\subsection{Gradient boosting i XGBoost}
W boostingu model jest sumą drzew budowanych kolejno, z których każde poprawia błędy poprzednich. XGBoost minimalizuje funkcję celu z regularyzacją \cite{chen2016}
\begin{equation}
  \mathcal{L}=\sum_i l(y_i,\hat{y}_i)+\sum_k \Omega(f_k),\qquad \Omega(f)=\gamma T+\tfrac{1}{2}\lambda\sum_{j=1}^{T}w_j^2,
\end{equation}
gdzie $T$ to liczba liści drzewa, a $w_j$ ich wagi. Istotne parametry to m.in. głębokość drzew, współczynnik uczenia i \texttt{scale\_pos\_weight}, służący do ważenia klasy rzadkiej.

\section{Algorytmy uczenia nienadzorowanego}
\subsection{Isolation Forest}
Metoda opiera się na obserwacji, że anomalie są „nieliczne i różne”, więc izoluje się je losowymi podziałami po mniejszej liczbie kroków niż punkty normalne \cite{liu2008}. Dla punktu $x$ i próby $n$ obserwacji wynik anomalii wynosi
\begin{equation}
  s(x,n)=2^{-\frac{E[h(x)]}{c(n)}},\qquad c(n)=2H(n-1)-\frac{2(n-1)}{n},
\end{equation}
gdzie $h(x)$ to długość ścieżki izolacji w drzewie, a $H(i)\approx\ln i+0{,}5772$ liczba harmoniczna. Wartości bliskie 1 wskazują anomalię. Złożoność uczenia jest liniowa względem liczby próbek.

\subsection{Local Outlier Factor}
LOF porównuje lokalną gęstość punktu z gęstością jego sąsiadów \cite{breunig2000}. Przy odległości osiągalności $\text{reach-dist}_k(p,o)=\max\{k\text{-dist}(o),d(p,o)\}$ lokalna gęstość to $\text{lrd}_k(p)=|N_k(p)|/\sum_{o\in N_k(p)}\text{reach-dist}_k(p,o)$, a
\begin{equation}
  \text{LOF}_k(p)=\frac{1}{|N_k(p)|}\sum_{o\in N_k(p)}\frac{\text{lrd}_k(o)}{\text{lrd}_k(p)}.
\end{equation}
Wartości istotnie większe od 1 oznaczają punkt w rzadszym otoczeniu niż jego sąsiedzi. Aby oceniać nowe obserwacje, stosuje się wariant wykrywania nowości (ang. \emph{novelty detection}).

\subsection{Metody oparte na rekonstrukcji i inne}
W badaniu \cite{niu2019} porównano także One-Class SVM, autoenkoder, ograniczoną maszynę Boltzmanna (RBM) i sieć GAN. Autoenkodery uczą się odtwarzać dane normalne, a wynikiem anomalii jest błąd rekonstrukcji. \todo{rozszerzyć o opis wybranych metod, jeśli wejdą do eksperymentów}

\section{Problem niezbalansowania klas}
Przy rzadkiej klasie dodatniej klasyfikator minimalizujący błąd globalny może ignorować oszustwa. Stosuje się: (i) ważenie klas, np. $w_c=n/(2n_c)$ dla klas binarnych, (ii) niedopróbkowanie klasy większościowej, (iii) nadpróbkowanie klasy rzadkiej, w tym SMOTE, który generuje syntetyczne obserwacje $\mathbf{x}_{\text{nowa}}=\mathbf{x}_i+\lambda(\mathbf{x}_{nn}-\mathbf{x}_i)$, $\lambda\sim U(0,1)$, na odcinku między obserwacją a jej sąsiadem z klasy rzadkiej \cite{chawla2002}. Próbkowanie stosuje się wyłącznie do zbioru treningowego. Zmiana proporcji klas zniekształca prawdopodobieństwa zwracane przez model, więc dobór progu wymaga ostrożności \cite{dalpozzolo2018}.

\section{Metryki ewaluacyjne dla klas niezbalansowanych}
Z macierzy pomyłek (TP, FP, FN, TN) wyznacza się
\begin{equation}
  \text{Precision}=\frac{TP}{TP+FP},\qquad \text{Recall}=\frac{TP}{TP+FN},\qquad F_1=\frac{2\,\text{Precision}\cdot\text{Recall}}{\text{Precision}+\text{Recall}}.
\end{equation}
Dokładność (ang. \emph{accuracy}) jest przy 0{,}17\% oszustw bezużyteczna, bo model oznaczający wszystkie transakcje jako normalne osiąga ponad 99{,}8\%. Krzywa ROC przedstawia czułość względem odsetka fałszywych alarmów, a ROC AUC jest polem pod nią. Krzywa precision--recall i średnia precyzja $\text{AP}=\sum_n(R_n-R_{n-1})P_n$ lepiej oddają jakość detekcji rzadkiej klasy, ponieważ nie uwzględniają prawdziwie ujemnych, a poziom losowy wynosi odsetek oszustw \cite{saito2015}. Próg decyzyjny wybiera się na zbiorze walidacyjnym, np. maksymalizując $F_1$ lub minimalizując koszt $C=c_{FN}\cdot FN+c_{FP}\cdot FP$.

\section{Przegląd istniejących rozwiązań (State-of-the-Art)}
Praca \cite{dalpozzolo2018}, powstała we współpracy z partnerem przemysłowym, formalizuje problem detekcji oszustw w warunkach rzeczywistych: niezbalansowanie, dryf koncepcji i opóźnienie weryfikacji oraz miary jakości dostosowane do pracy z alertami. Strategię uczenia przetestowano na ponad 75 mln transakcji e-commerce.

Praca \cite{niu2019} porównuje na zbiorze \cite{kaggle_ulb} sześć modeli nadzorowanych i cztery nienadzorowane, stosując miarę AUROC. Najlepsze wyniki uzyskały XGBoost i Random Forest, a wśród modeli nienadzorowanych RBM (tab.~\ref{tab:niu}). Autorzy wnioskują, że metody nadzorowane są nieznacznie lepsze, lecz nienadzorowane są obiecujące, gdy brakuje etykiet. Zastosowano losowe niedopróbkowanie i losową walidację krzyżową.

\begin{table}[H]
\centering
\caption{Wybrane wyniki (AUROC) z pracy \cite{niu2019} dla zbioru Kaggle.}
\label{tab:niu}
\begin{tabular}{lll}
\toprule
Model & Rodzaj & AUROC \\
\midrule
XGBoost & nadzorowany & 0{,}989 \\
Random Forest & nadzorowany & 0{,}988 \\
RBM & nienadzorowany & 0{,}961 \\
\bottomrule
\end{tabular}\\[4pt]
\zrodlo{\cite{niu2019}}
\end{table}

Praca \cite{carcillo2019} pokazuje, że łączenie technik nadzorowanych i nienadzorowanych (wyniki odstępstwa liczone na różnych poziomach granulacji) poprawia skuteczność wykrywania na rzeczywistym, oznaczonym zbiorze.

\todo{dodać 1--2 nowsze prace przeglądowe (2022 i później) oraz wskazać lukę badawczą.}

\paragraph{Miejsce niniejszej pracy.} W odróżnieniu od części publikacji porównawczych, które ograniczają się do ewaluacji offline z losowym podziałem i miarą ROC AUC \cite{niu2019}, w pracy zastosowano podział czasowy, PR-AUC, wybór progu na walidacji i analizę niepewności wyników, a porównanie uzupełniono działającym systemem webowym.

\chapter{Projekt systemu i metodyka badawcza}
\label{ch:projekt}

\section{Założenia i wymagania systemowe}
\todo{Wymagania funkcjonalne (ocena transakcji, zarządzanie modelami, metryki, historia) i niefunkcjonalne (opóźnienie, modułowość); tabela.}

\section{Architektura klient--serwer i dobór narzędzi}
\todo{Diagram: frontend (React), API (FastAPI), moduł ML (scikit-learn, XGBoost), baza (SQL).}
Architektura jest modułowa: źródła transakcji (symulator, odtwarzanie zbioru, import CSV, szkielet konektora bankowego) dostarczają ujednoliconych obiektów \texttt{Transaction}, a moduł oceny mapuje je na cechy właściwe dla schematu danego modelu. Dzięki temu dodanie nowego źródła lub modelu nie wymaga zmian pozostałych części. \todo{rozwinąć po ustabilizowaniu kodu}

\section{Model relacyjnej bazy danych}
\todo{Schemat tabel (historia ocen, modele) i uzasadnienie SQLite/PostgreSQL.}

\section{Charakterystyka i podział zbioru danych}
Zbiór \emph{Credit Card Fraud Detection} \cite{kaggle_ulb} zawiera 284\,807 transakcji kartowych z dwóch dni, w tym 492 oszustwa (0{,}172\%). Cechy V1--V28 są wynikiem transformacji PCA (ze względu na poufność), a ponadto dostępne są czas od pierwszej transakcji, kwota i etykieta klasy. \todo{cytowanie zbioru zgodnie z opisem na stronie, liczność klas w podzbiorach, uwagi o ograniczeniach.}

Dane dzielono chronologicznie na zbiory treningowy, walidacyjny i testowy (70/15/15), tak aby testowanie odpowiadało predykcji „przyszłości” względem uczenia. \todo{rozkład oszustw w podzbiorach, opis wariantu odtwarzającego ustawienie z \cite{niu2019}}

\chapter{Implementacja systemu}
\label{ch:impl}

\section{Moduł uczenia maszynowego i zarządzania modelami}
\todo{Rejestr modeli, trening, zapis metryk i progu, wersjonowanie.}

\section{Backend i interfejs API}
\todo{Tabela endpointów (predykcja, modele, metryki, sandbox, strumień).}

\section{Panel użytkownika}
\todo{Widoki: modele, metryki, predykcja, dashboard transakcji.}

\section{Scenariusz testów i weryfikacja poprawności}
\todo{Testy jednostkowe i integracyjne, pomiar opóźnienia.}

\chapter{Badania własne i analiza wyników}
\label{ch:badania}

% Tabele i wykresy generuje: python -m anomaly.experiments (skopiuj results/tables -> tables/, results/figures -> figures/)
% \input{tables/tab_podzial.tex}   \input{tables/tab_wyniki_glowne.tex}   \input{tables/tab_przedzialy.tex}
% \input{tables/tab_niu.tex}   \input{tables/tab_niezbalansowanie.tex}   \input{tables/tab_okna.tex}
\section{Metodologia eksperymentu}
\todo{Podział danych, walidacja, hiperparametry, powtórzenia (seedy), sprzęt.}
\section{Skuteczność algorytmów uczenia nadzorowanego}
\todo{Tabela metryk, macierze pomyłek.}
\section{Skuteczność algorytmów uczenia nienadzorowanego}
\todo{Tabela metryk.}
\section{Wpływ sposobu radzenia sobie z niezbalansowaniem klas}
\todo{Wagi klas vs SMOTE.}
\section{Porównanie wyników}
\todo{Krzywe ROC/PR, wpływ progu, porównanie z \cite{niu2019}.}
\section{Analiza statystyczna wyników}
\todo{Powtórzenia i bootstrap; przedziały ufności.}
\section{Ocena wydajności systemu i ograniczenia}
\todo{Opóźnienie oceny, ograniczenia zbioru (cechy PCA, dwa dni, brak dryfu).}

\chapter*{Podsumowanie}
\addcontentsline{toc}{chapter}{Podsumowanie}

\todo{Co zrobiono i jakie rezultaty osiągnięto (własnymi słowami, na końcu pracy).}

\todo{Odpowiedzi na pytania badawcze zadane we wstępie.}

\todo{Kierunki dalszego rozwoju: sandbox z własnymi transakcjami, konektor bankowy (Open Banking), podejście hybrydowe \cite{carcillo2019}, dryf koncepcji.}


\printbibliography[heading=bibintoc,title={Bibliografia}]

\chapter*{Załącznik A}
\addcontentsline{toc}{chapter}{Załącznik A}

Spis zawartości dołączonego nośnika (repozytorium kodu):
\begin{itemize}
  \item tekst pracy w formacie PDF,
  \item \texttt{anomaly/} -- moduł uczenia maszynowego (dane, modele, metryki, trening),
  \item \texttt{backend/} -- interfejs API (FastAPI),
  \item \texttt{notebooks/} -- eksperymenty i analiza danych,
  \item \texttt{tests/} -- testy automatyczne.
\end{itemize}

Komendy przydatne podczas uruchamiania programu:
\begin{lstlisting}
pip install -r requirements.txt     # instalacja środowiska
python -m anomaly.train --data data/raw/creditcard.csv --models all
uvicorn backend.app.main:app --reload     # uruchomienie API
pytest                                     # testy
\end{lstlisting}
\todo{uzupełnić po ukończeniu implementacji (frontend, dane)}

\end{document}
